\section*{Chapter \ref{ch:g11:trigo} --- Trigonometry: The Unit Circle}

\begin{solution}{exo:g11:trigo:1}
To \omterm{def:g11:trigo:radian}{radians}: $30^\circ = \frac{\pi}{6}$, $45^\circ = \frac{\pi}{4}$,
$120^\circ = \frac{2\pi}{3}$, $270^\circ = \frac{3\pi}{2}$.

To degrees: $\frac{\pi}{5} = 36^\circ$, $\frac{5\pi}{6} = 150^\circ$,
$\frac{7\pi}{4} = 315^\circ$, $\frac{2\pi}{9} = 40^\circ$.
\end{solution}

\begin{solution}{exo:g11:trigo:2}
$\frac{2\pi}{3}$ and $-\frac{\pi}{4}$ are already in
$\intoc{-\pi}{\pi}$: second quadrant and fourth quadrant respectively.

$\frac{17\pi}{6} - 2\pi = \frac{5\pi}{6}$: second quadrant, close to the
negative $x$-axis.

$-\frac{7\pi}{2} + 4\pi = \frac{\pi}{2}$: the top of the circle,
$(0, 1)$.
\end{solution}

\begin{solution}{exo:g11:trigo:3}
$\cos\frac{3\pi}{4} = \cos\left(\pi - \frac\pi4\right)
= -\cos\frac\pi4 = -\frac{\sqrt2}{2}$.

$\sin\frac{5\pi}{6} = \sin\left(\pi - \frac\pi6\right)
= \sin\frac\pi6 = \frac12$.

$\cos\left(-\frac\pi3\right) = \cos\frac\pi3 = \frac12$.

$\sin\frac{4\pi}{3} = \sin\left(\pi + \frac\pi3\right)
= -\sin\frac\pi3 = -\frac{\sqrt3}{2}$.

$\cos\frac{11\pi}{6} = \cos\left(-\frac\pi6 + 2\pi\right)
= \cos\frac\pi6 = \frac{\sqrt3}{2}$.
\end{solution}

\begin{solution}{exo:g11:trigo:4}
$\sin^2 t = 1 - \cos^2 t = 1 - \frac{25}{169} = \frac{144}{169}$, so
$\sin t = \pm\frac{12}{13}$. On $\intoo{-\frac\pi2}{0}$ (fourth quadrant)
the \omterm{def:g10:coordgeom:system}{ordinate} is negative: $\sin t = -\frac{12}{13}$.
\end{solution}

\begin{solution}{exo:g11:trigo:5}
Using \cref{prop:g11:trigo:associated}:
\[
A = (-\cos x) + \cos x + \cos x - \cos x = 0 ,
\]
\[
B = \sin x + (-\sin x) + (-\sin x) = -\sin x .
\]
\end{solution}

\begin{solution}{exo:g11:trigo:6}
$\cos t = \frac{\sqrt3}{2}$: one solution is $\frac\pi6$, so
$t = \pm\frac\pi6$ (the shifts by $2k\pi$ leave the \omterm{def:g10:numbers:interval}{interval}).

$\sin t = \frac12$: $t = \frac\pi6$ or $t = \pi - \frac\pi6 =
\frac{5\pi}{6}$.

$\cos t = -1$: the only point of \omterm{def:g10:coordgeom:system}{abscissa} $-1$ is $M(\pi)$, so $t = \pi$.
\end{solution}

\begin{solution}{exo:g11:trigo:7}
$\sin t = -\frac{\sqrt3}{2}$: from $a = -\frac\pi3$, the families are
$-\frac\pi3 + 2k\pi$ and $\pi + \frac\pi3 + 2k\pi$; in
$\intco{0}{2\pi}$: $t = \frac{4\pi}{3}$ and $t = \frac{5\pi}{3}$.

$\cos t = 0$: $t = \frac\pi2$ and $t = \frac{3\pi}{2}$.

$2\sin t + 1 = 0$ means $\sin t = -\frac12$: from $a = -\frac\pi6$, in
$\intco{0}{2\pi}$: $t = \frac{7\pi}{6}$ and $t = \frac{11\pi}{6}$.
\end{solution}

\begin{solution}{exo:g11:trigo:8}
By \cref{thm:g11:trigo:equations}, $\cos 2t = \cos t$ means
$2t = t + 2k\pi$ or $2t = -t + 2k\pi$, i.e.\ $t = 2k\pi$ or
$t = \frac{2k\pi}{3}$ ($k \in \Z$). The first family is contained in the
second (take $k$ a multiple of $3$), so the solution set is
\[
S = \left\{\frac{2k\pi}{3} : k \in \Z\right\}
\]
--- on the circle, the three points $M(0)$,
$M\!\left(\frac{2\pi}{3}\right)$, $M\!\left(\frac{4\pi}{3}\right)$.
\end{solution}

\begin{solution}{exo:g11:trigo:9}
\omterm{def:g10:algebra:expand}{Expanding} both squares and using $\cos^2 x + \sin^2 x = 1$:
\[
(\cos^2 x + 2\sin x\cos x + \sin^2 x)
+ (\cos^2 x - 2\sin x\cos x + \sin^2 x)
= 1 + 1 = 2 .
\]
\end{solution}

\begin{solution}{exo:g11:trigo:10}
Rolling without slipping means the arc length equals the distance
travelled: $r\theta = 150$ cm gives
$\theta = \frac{150}{30} = 5$ rad $= 5 \times \frac{180}{\pi} \approx
286.5^\circ$. One full turn advances the bike by the circumference
$2\pi \times 30 = 60\pi \approx 188.5$ cm, about $1.88$ m.
\end{solution}

\begin{solution}{exo:g11:trigo:11}
On the unit circle, the points of \omterm{def:g10:coordgeom:system}{abscissa} exactly $\frac12$ are
$M\!\left(\pm\frac\pi3\right)$. The points with \omterm{def:g10:coordgeom:system}{abscissa} \emph{at most}
$\frac12$ form the arc on the left of the vertical line $x = \frac12$,
travelled as $t$ goes from $\frac\pi3$ up to $\pi$, or from $-\pi$ up to
$-\frac\pi3$. Hence, on $\intoc{-\pi}{\pi}$:
\[
S = \intoc{-\pi}{-\frac{\pi}{3}} \cup \intcc{\frac{\pi}{3}}{\pi}.
\]
(The endpoints $\pm\frac\pi3$ are included since the inequality is wide.)
\end{solution}

\begin{solution}{pb:g11:trigo:1}
\textbf{1.} $30^\circ = \frac{\pi}{6}$; $135^\circ =
\frac{3\pi}{4}$; $300^\circ = \frac{5\pi}{3}$. And
$\frac{3\pi}{4} = 135^\circ$; $\frac{5\pi}{6} = 150^\circ$.

\textbf{2.} Length $r\,t = 3 \times 2 = 6$ m. \omterm{def:g11:trigo:radian}{Radians} make arc
length a plain product --- degrees would drag a
$\frac{\pi}{180}$ into every formula.

\textbf{3.} \omterm{def:g11:trigo:radian}{Radians} turned $= \frac{\text{distance}}
{\text{radius}} = \frac{1000}{0.30} \approx 3\,333$ rad;
divided by $2\pi$: about $530$ full turns.

\textbf{4.} $\cos\frac{2\pi}{3} = -\frac12$;
$\sin\left(-\frac\pi4\right) = -\frac{\sqrt2}{2}$;
$\frac{19\pi}{6} - 2\pi = \frac{7\pi}{6}$, so
$\cos\frac{19\pi}{6} = \cos\frac{7\pi}{6} =
-\frac{\sqrt3}{2}$.

\textbf{5.} $\cos^2 t = 1 - \frac{9}{25} = \frac{16}{25}$; in
the second quadrant the \omterm{def:g11:trigo:cossin}{cosine} is negative:
$\cos t = -\frac45$, and
$\tan t = \frac{3/5}{-4/5} = -\frac34$.

\textbf{6.} In $t$ minutes the wheel turns
$\frac{t}{30} \times 2\pi = \frac{\pi t}{15}$. At $t = 0$:
$h = 65 - 60 = 5$ m --- the boarding platform at the bottom
(hub height minus radius). The minus sign puts you \emph{below}
the hub when the turned angle is $0$.

\textbf{7.} $h(7.5) = 65 - 60\cos\frac{\pi}{2} = 65$ m (hub
height); $h(15) = 65 + 60 = 125$ m (the top);
$h(20) = 65 - 60\cos\frac{4\pi}{3} = 65 + 30 = 95$ m.

\textbf{8.} $65 - 60\cos\frac{\pi t}{15} = 95$ gives
$\cos\frac{\pi t}{15} = -\frac12$: first solution
$\frac{\pi t}{15} = \frac{2\pi}{3}$, i.e.\ $t = 10$ minutes.

\textbf{9.} $\cos\frac{\pi t}{15} \leq -\frac12$ for
$\frac{\pi t}{15} \in \intcc{\frac{2\pi}{3}}{\frac{4\pi}{3}}$,
i.e.\ $t \in \intcc{10}{20}$: ten minutes of the ride above
$95$ m, symmetric about the summit at $t = 15$.

\textbf{10.} First minute:
\[
h(1) - h(0) = 60\left(1 - \cos\frac{\pi}{15}\right)
\approx 1.3 \text{ m}.
\]
Between minutes $7$ and $8$:
\[
60\left(\cos\frac{7\pi}{15} - \cos\frac{8\pi}{15}\right)
\approx 12.5 \text{ m}
\]
--- nearly ten times more. The height changes fastest when passing hub level
and stalls near bottom and top: the sinusoid is steep at its
middle, flat at its extremes.

\textbf{11.} Radius $= \frac{40 - 4}{2} = 18$ m, hub at
$4 + 18 = 22$ m, period $12$ min:
$h(t) = 22 - 18\cos\left(\frac{\pi t}{6}\right)$.

\textbf{12.} \omterm{def:g10:functions:extrema}{Maximum} $7 + 3 = 10$ m when
$\sin\frac{\pi t}{6} = 1$: first at $t = 3$ (3 a.m.);
\omterm{def:g10:functions:extrema}{minimum} $4$ m first at $t = 9$; period
$\frac{2\pi}{\pi/6} = 12$ hours.

\textbf{13.} $\sin\frac{\pi t}{6} \geq \frac12$ for
$\frac{\pi t}{6} \in \intcc{\frac\pi6}{\frac{5\pi}{6}}$:
$t \in \intcc{1}{5}$ --- a four-hour window from 1 a.m. to 5
a.m., reopening one period later, $t \in \intcc{13}{17}$
(1 p.m. to 5 p.m.).

\textbf{14.} Mirror across the vertical axis sends the point at
angle $t$ to the point at angle $\pi - t$, preserving the
\omterm{def:g10:coordgeom:system}{ordinate}: $\sin(\pi - t) = \sin t$. The half-turn about the
center sends $t$ to $\pi + t$, negating both \omterm{def:g10:coordgeom:system}{coordinates}:
$\cos(\pi + t) = -\cos t$. And the mirror across the diagonal
swaps \omterm{def:g10:coordgeom:system}{abscissa} and \omterm{def:g10:coordgeom:system}{ordinate}:
$\cos\left(\frac\pi2 - t\right) = \sin t$ --- the
``co'' of complementary angles.

\textbf{15.} $\sin(-t) = -\sin t$: odd (half-turn symmetry of
its \omterm{def:g10:functions:graph}{graph}); $\cos(-t) = \cos t$: even (mirror). And
$\sin t + t^3$ is a sum of two \omterm{def:g11:func:parity}{odd functions}: odd.

\textbf{16.} Factor: $(2\sin t + 1)(\sin t - 1) = 0$:
$\sin t = 1$ gives $t = \frac\pi2$;
$\sin t = -\frac12$ gives $t = \frac{7\pi}{6}$ and
$t = \frac{11\pi}{6}$. Three solutions on the circle.

\textbf{17.} $1$ rad $= \frac{180}{\pi} \approx 57.3^\circ$.
So $\sin(1) \approx 0.841$ while $\sin(1^\circ) \approx
0.0175$: a factor of nearly $50$. Moral: check the mode light
before trusting any trigonometric display.

\textbf{18.} \omterm{def:g10:coordgeom:system}{Abscissa} equals \omterm{def:g10:coordgeom:system}{ordinate} on the diagonal:
$t = \frac\pi4$ and $t = \frac{5\pi}{4}$.

\textbf{19.} $\sin(0.1) = 0.09983\ldots$ against $0.1$:
relative error about $0.17\,\%$. For small angles the chord
hugs the arc --- and the grade-12 limit
$\frac{\sin t}{t} \to 1$ is exactly this observation made into
a theorem.

\textbf{20.} \omterm{def:g11:trigo:radian}{Radians}: angle $=$ arc on the unit circle, so
angles compute like lengths. \omterm{def:g11:trigo:cossin}{Cosine} and \omterm{def:g11:trigo:cossin}{sine}: the \omterm{def:g10:coordgeom:system}{coordinates}
of the clock hand --- shadow on the ground, height on the
wall. Symmetries of the circle: the whole catalogue of
identities, read off mirrors and half-turns. \omterm{def:g10:algebra:equation}{Equations}
$\cos t = c$: the clock's timetable, with the circle showing
every solution at once. Costumes: the wheel wore
$65 - 60\cos$, the tide wore $7 + 3\sin$ --- same clock, same
mathematics.
\end{solution}
